% S3. Order book de BTC/USDT: dos lados, niveles de precio y qué es un delta.
\begin{tikzpicture}[x=1cm,y=1cm]
  \path[use as bounding box] (0,-0.1) rectangle (13.5,4.75);
  % ---------- Libro ----------
  \node[anchor=west,font=\scriptsize\bfseries,text=colVenta] at (0.6,4.5) {VENTAS (\emph{asks})};
  \node[venta] (a2) at (2.3,3.85) {67\,002 USD \quad 0.5 BTC};
  \node[venta] (a1) at (2.3,3.25) {67\,001 USD \quad 2.0 BTC};
  \draw[colNeutro,dashed,thick] (0.4,2.75) -- (4.2,2.75);
  \node[compra] (b1) at (2.3,2.25) {67\,000 USD \quad 1.2 BTC};
  \node[compra] (b2) at (2.3,1.65) {66\,999 USD \quad 0.8 BTC};
  \node[anchor=west,font=\scriptsize\bfseries,text=colCompra] at (0.6,1.05) {COMPRAS (\emph{bids})};
  % Precio de mercado
  \node[font=\scriptsize,text=colAcento,align=left,anchor=west] (pm) at (4.45,2.75) {precio de mercado\\entre 67\,000 y 67\,001};
  % Nivel = suma de órdenes
  \node[font=\scriptsize,align=left,anchor=north west,text width=6.5cm] at (0.1,0.7) {Cada fila es un \textbf{nivel de precio} y suma todas las órdenes a ese precio, por ejemplo 2.0 = 0.5 + 1.5 de dos vendedores.};
  % ---------- Qué cambia el libro ----------
  \draw[colNeutro!40] (7.0,-0.05) -- (7.0,4.7);
  \node[anchor=west,font=\small\bfseries] at (7.3,4.5) {Qué cambia el libro};
  \node[box,font=\scriptsize,text width=5.4cm] (s1) at (10.35,3.75) {1. Llega una orden de compra de \textbf{1 BTC} al mejor precio};
  \node[box,font=\scriptsize,text width=5.4cm] (s2) at (10.35,2.75) {2. La orden se cruza con el vendedor más barato en \textbf{67\,001}};
  \node[box,font=\scriptsize,text width=5.4cm] (s3) at (10.35,1.75) {3. El nivel 67\,001 pasa de 2.0 a \textbf{1.0 BTC}};
  \node[acento,font=\scriptsize,text width=5.4cm] (s4) at (10.35,0.6) {Ese cambio es un \textbf{delta}, que trae el valor nuevo\\\emph{``67\,001 ahora vale 1.0''}};
  \draw[flow] (s1) -- (s2);
  \draw[flow] (s2) -- (s3);
  \draw[flow] (s3) -- (s4);
\end{tikzpicture}
